% =====================================================================
%  MASTER LaTeX EXAM CHEAT SHEET
%  A single compilable reference covering every feature likely to appear
%  in a fast, 50-minute typesetting test. Read it, compile it, keep the
%  PDF open beside you during the exam.
%
%  Compile: pdflatex cheatsheet.tex  (run twice for TOC / refs)
%  Some sections need: -shell-escape (only if you switch listings->minted)
% =====================================================================
\documentclass[11pt]{article}

% ---------------- CORE PREAMBLE (copy this every time) ----------------
\usepackage[a4paper,margin=1in]{geometry}
\usepackage[T1]{fontenc}
\usepackage[utf8]{inputenc}
\usepackage{amsmath,amssymb,amsfonts,amsthm}
\usepackage{graphicx}
\usepackage{xcolor}
\usepackage{enumitem}
\usepackage{hyperref}          % keep hyperref near-last

% ---------------- FEATURE PACKAGES ----------------
\usepackage{subcaption}        % subfigures
\usepackage{wrapfig}           % text-wrapped figures
\usepackage{rotating}          % rotated floats
\usepackage{array}             % column tuning
\usepackage{multirow}          % row spanning in tables
\usepackage{booktabs}          % \toprule \midrule \bottomrule
\usepackage{tabularx}          % auto-width columns
\usepackage{colortbl}          % \rowcolor \cellcolor \columncolor
\usepackage{listings}          % code listings
\usepackage{authblk}           % numbered author affiliations
\usepackage{tcolorbox}         % coloured boxes / custom theorem boxes
\usepackage[ruled,vlined,linesnumbered]{algorithm2e} % pseudocode

\hypersetup{colorlinks=true, linkcolor=blue, urlcolor=magenta, citecolor=teal}

% ---------------- CUSTOM COMMANDS (save keystrokes) ----------------
\newcommand{\R}{\mathbb{R}}            % \R -> real numbers
\newcommand{\dd}{\,\mathrm{d}}         % upright differential
\newcommand{\pd}[2]{\frac{\partial #1}{\partial #2}} % partial derivative
\DeclareMathOperator{\Lap}{\mathcal{L}} % a named operator

% ---------------- THEOREM ENVIRONMENTS ----------------
\theoremstyle{plain}
\newtheorem{theorem}{Theorem}[section]
\newtheorem{proposition}[theorem]{Proposition}
\newtheorem{lemma}[theorem]{Lemma}
\newtheorem{corollary}[theorem]{Corollary}
\theoremstyle{definition}
\newtheorem{definition}{Definition}[section]
\theoremstyle{remark}
\newtheorem*{remark}{Remark}

% A coloured tcolorbox theorem for "custom theorem box" questions
\newtcolorbox{keybox}[1][]{colback=yellow!8, colframe=orange!70!black,
  fonttitle=\bfseries, title=#1}

% ---------------- LISTINGS STYLE ----------------
\lstdefinestyle{code}{
  basicstyle=\ttfamily\small, keywordstyle=\color{blue}\bfseries,
  commentstyle=\color{green!50!black}\itshape, stringstyle=\color{red!70!black},
  numbers=left, numberstyle=\tiny\color{gray}, numbersep=8pt,
  frame=single, breaklines=true, showstringspaces=false, tabsize=2}

% =====================================================================
\title{\textbf{Master \LaTeX{} Exam Cheat Sheet}}
\author{Prepared for a 50-minute typesetting test}
\date{\today}
% =====================================================================

\begin{document}
\maketitle
\tableofcontents
\newpage

% =====================================================================
\section{Fast Preamble Templates}

% --- Minimal (article, math + graphics) ---
% \documentclass[11pt]{article}
% \usepackage[margin=1in]{geometry}
% \usepackage{amsmath,amssymb,graphicx,xcolor,hyperref}
% \begin{document} ... \end{document}

% --- Exam-standard (everything you usually need) ---
% see the CORE + FEATURE preamble at the top of this file.

The rule of thumb: load \texttt{amsmath} for equations,
\texttt{graphicx} for images, \texttt{xcolor} for colour,
\texttt{hyperref} LAST. Add \texttt{multirow}, \texttt{subcaption},
\texttt{listings} only when the question needs them.

% =====================================================================
\section{Document Structure}
\subsection{Title block}
Use either \texttt{\textbackslash author} with \texttt{\textbackslash thanks}
or the \texttt{authblk} package for numbered affiliations:
\begin{verbatim}
\title{My Title}
\author[1]{Alice}\author[2]{Bob}
\affil[1]{University A}
\affil[2]{University B}
\date{}            % remove date
\maketitle
\end{verbatim}

\subsection{Sectioning \& Table of Contents}
\verb|\section \subsection \subsubsection \paragraph|. Star versions
(\verb|\section*|) are unnumbered and skip the TOC. \verb|\tableofcontents|
prints a clickable TOC when \texttt{hyperref} is loaded.

% =====================================================================
\section{Text Formatting \& Colour}
\textbf{bold}, \textit{italic}, \underline{underline},
\texttt{monospace}, \textsc{Small Caps}, \emph{emphasis},
{\Large large} and {\tiny tiny} sizes.
Colour: {\color{red}red text}, \textcolor{blue}{blue text},
\colorbox{yellow}{highlighted}, \fcolorbox{black}{green!20}{boxed}.
Define your own: \verb|\definecolor{myblue}{RGB}{0,80,180}|.

Size ladder (small $\to$ big): \verb|\tiny \scriptsize \footnotesize|
\verb|\small \normalsize \large \Large \LARGE \huge \Huge|.

% =====================================================================
\section{Lists \& Nested Lists}
\begin{enumerate}[label=(\alph*)]     % (a) (b) (c)
  \item Ordered with custom label.
  \item Second.
    \begin{itemize}[label=$\star$]    % star bullets
      \item bullet with star.
      \item[$\oplus$] one-off different marker via \verb|\item[...]|.
    \end{itemize}
\end{enumerate}
\begin{enumerate}[label=\Roman*.]     % I. II. III.
  \item Roman numerals.
\end{enumerate}
\begin{description}
  \item[Term] definition-style list (bold term + text).
\end{description}
Handy \texttt{enumitem} keys: \verb|label=|, \verb|leftmargin=|,
\verb|itemsep=2pt|, \verb|noitemsep|, \verb|resume| (continue numbering),
\verb|start=3| (begin at 3).

% =====================================================================
\section{Mathematics}
\subsection{Modes}
Inline math with \verb|$...$|: $E=mc^2$.
Display with \verb|\[ ... \]| (never use \verb|$$...$$| in modern LaTeX):
\[ \int_0^\infty e^{-x}\dd x = 1. \]
Numbered single equation with \texttt{equation}; multi-line with
\texttt{align} (\verb|&| = alignment point, \verb|\\| = new line).

\subsection{Alignment environments}
\begin{align}
  f(x) &= (x+1)^2            \label{eq:cs1}\\
       &= x^2 + 2x + 1 .     \label{eq:cs2}
\end{align}
\verb|align*| suppresses all numbers; \verb|\nonumber| suppresses one line;
\verb|\tag{$\ast$}| forces a custom tag. \texttt{gather} centres each line
with its own number; \texttt{multline} splits one long equation.
\[
  \begin{gathered}         % grouped, single number outside
    a = b + c \\
    d = e - f
  \end{gathered}
\]

\subsection{Fractions, roots, operators}
$\dfrac{a}{b}$ (display size) vs $\tfrac{a}{b}$ (text size),
$\sqrt{x}$, $\sqrt[n]{x}$, $\displaystyle\sum_{i=1}^{n} i$,
$\displaystyle\prod_{k}$, $\displaystyle\lim_{x\to\infty}$,
$\displaystyle\int_a^b$, $\oint$, $\iint$, $\binom{n}{k}$.

\subsection{Sub/superscripts, deep nesting}
Group with braces: \verb|x^{2n+1}|, \verb|a_{i,j}|. Nest freely:
$x_{i_{j_{k}}}$, $e^{-t/(RC)}$, $A^{\mu}_{\ \nu}$ (aligned tensor indices
via a phantom space \verb|\ |), $T^{2}_{M}$. Multi-line subscript:
$\displaystyle\sum_{\substack{i=1\\ i\ne j}}^{n}$.

\subsection{Delimiters that grow}
\verb|\left( \right)| auto-size: $\left(\dfrac{a}{b}\right)$,
$\left[\;\right]$, $\left\{\;\right\}$, $\left|\;\right|$,
$\left\langle\;\right\rangle$, $\left\lfloor x \right\rfloor$,
$\left\lceil x \right\rceil$. Manual sizes: \verb|\big \Big \bigg \Bigg|.

\subsection{Matrices \& cases}
\[
  \begin{pmatrix} a & b \\ c & d \end{pmatrix}
  \begin{bmatrix} 1 & 0 \\ 0 & 1 \end{bmatrix}
  \begin{vmatrix} a & b \\ c & d \end{vmatrix}
  \quad
  f(x)=\begin{cases} x^2 & x\ge 0\\ -x & x<0 \end{cases}
\]
\texttt{pmatrix}=( ), \texttt{bmatrix}=[ ], \texttt{Bmatrix}=\{ \},
\texttt{vmatrix}=| |, \texttt{Vmatrix}=$\|\ \|$, \texttt{matrix}=none.

\subsection{Spacing in math}
Thin \verb|\,| medium \verb|\:| thick \verb|\;| negative \verb|\!| quad
\verb|\quad| \verb|\qquad|. Text inside math: \verb|\text{...}|.
Upright differential: \verb|\mathrm{d}|.

\subsection{Symbol quick reference}
Greek: $\alpha\beta\gamma\delta\epsilon\varepsilon\zeta\eta\theta
\vartheta\iota\kappa\lambda\mu\nu\xi\pi\varpi\rho\varrho\sigma\varsigma
\tau\upsilon\phi\varphi\chi\psi\omega$ and caps
$\Gamma\Delta\Theta\Lambda\Xi\Pi\Sigma\Upsilon\Phi\Psi\Omega$.\\
Relations: $\le\ge\ne\approx\equiv\sim\simeq\cong\propto\ll\gg
\subset\subseteq\supset\supseteq\in\notin\ni\prec\succ\preceq\succeq$.\\
Arrows: $\to\gets\Rightarrow\Leftarrow\Leftrightarrow\iff\implies
\mapsto\longrightarrow\xrightarrow{f}\uparrow\downarrow\rightleftharpoons$.\\
Operators/sets: $\cup\cap\uplus\setminus\times\div\pm\mp\cdot\ast\star
\oplus\otimes\odot\wedge\vee\forall\exists\nexists\emptyset\varnothing
\nabla\partial\infty\angle\triangleright\triangleleft\diamond\dagger$.\\
Blackboard/script/frak: $\mathbb{R}\ \mathcal{L}\ \mathfrak{g}$
(need \texttt{amssymb}/\texttt{amsfonts}). Dots: $\dots\cdots\vdots\ddots$.
Decorations: $\hat{a}\ \bar{a}\ \tilde{a}\ \vec{a}\ \dot{a}\ \widehat{abc}
\ \overline{abc}\ \underline{abc}\ \overrightarrow{abc}$.

% =====================================================================
\section{Tables}
\subsection{Basic \texttt{tabular}}
\begin{center}
\begin{tabular}{|l|c|r|}     % left, center, right; | = vertical rule
  \hline
  Left & Centre & Right \\ \hline
  a & b & c \\ \hline
\end{tabular}
\end{center}
\verb|&| separates columns, \verb|\\| ends a row, \verb|\hline| draws a
full rule, \verb|\cline{2-3}| a partial rule over columns 2--3.

\subsection{\texttt{multirow} + \texttt{multicolumn} + colour}
\begin{center}
\renewcommand{\arraystretch}{1.4}
\begin{tabular}{|c|c|c|}
  \hline
  \rowcolor{gray!30}
  \multicolumn{2}{|c|}{\textbf{Spanning header}} & \textbf{C} \\ \hline
  \multirow{2}{*}{Merged}
    & \cellcolor{yellow!30}$x^2$ & $s>0$ \\ \cline{2-3}
    & $\dfrac{1}{s}$             & $s>0$ \\ \hline
\end{tabular}
\end{center}
\verb|\multirow{N}{*}{txt}| spans N rows (\verb|*| = natural width);
\verb|\multicolumn{N}{|c|}{txt}| spans N columns and re-declares the
border; \verb|\rowcolor|, \verb|\cellcolor|, \verb|\columncolor| shade.
Math goes in cells simply by using \verb|$...$| or \verb|\dfrac|.

\subsection{Line thickness \& widths}
Thicker rules: \verb|\setlength{\arrayrulewidth}{1pt}| or use
\texttt{booktabs} (\verb|\toprule \midrule \bottomrule|).
Auto-fit to page: \texttt{tabularx} with an \texttt{X} column:
\begin{center}
\begin{tabularx}{0.8\textwidth}{|l|X|}
  \hline Key & Long wrapping description column \\ \hline
\end{tabularx}
\end{center}

% =====================================================================
\section{Figures \& Graphics}
\subsection{Single figure}
\begin{verbatim}
\begin{figure}[htbp]
  \centering
  \includegraphics[width=0.6\textwidth]{example-image-a}
  \caption{A caption.}\label{fig:one}
\end{figure}
\end{verbatim}
Scaling keys: \verb|width=|, \verb|height=|, \verb|scale=|,
\verb|angle=45| (rotate), \verb|keepaspectratio|. Float placement
\verb|[htbp]| = here/top/bottom/page; add \verb|!| to override LaTeX's
reluctance, or \verb|[H]| (needs \texttt{float}) to force exact spot.

\subsection{Side-by-side \& subfigures}
\begin{figure}[htbp]
  \centering
  \begin{subfigure}[b]{0.3\textwidth}
    \includegraphics[width=\linewidth]{example-image-a}
    \caption{First}\label{fig:sa}
  \end{subfigure}\hfill
  \begin{subfigure}[b]{0.3\textwidth}
    \includegraphics[width=\linewidth]{example-image-b}
    \caption{Second}\label{fig:sb}
  \end{subfigure}
  \caption{Two subfigures.}\label{fig:subs}
\end{figure}
\verb|\ref{fig:sa}| $\to$ ``\ref{fig:sa}''; \verb|\subref{fig:sa}| gives
just the sub-letter. \texttt{minipage} works too when you do not want
sub-captions:
\verb|\begin{minipage}{0.48\textwidth}...\end{minipage}|.

\subsection{Wrapped \& rotated figures}
\begin{verbatim}
\begin{wrapfigure}{r}{0.4\textwidth}   % r=right, l=left
  \centering
  \includegraphics[width=\linewidth]{example-image}
  \caption{Wrapped.}\label{fig:wrap}
\end{wrapfigure}
% rotated:
\includegraphics[angle=90]{example-image}   % rotate the image
\begin{sidewaysfigure}...\end{sidewaysfigure} % rotate whole float (rotating)
\end{verbatim}

% =====================================================================
\section{Cross-References \& Hyperlinks}
Label anything with \verb|\label{type:name}| directly after its
\verb|\caption| or in the equation, then cite with \verb|\ref| (number),
\verb|\eqref| (adds parentheses for equations), \verb|\pageref| (page),
or \verb|\autoref| (adds the word ``Figure''/``Table'').
External link: \verb|\href{https://url}{clickable text}|; bare URL:
\verb|\url{https://url}|. Configure colours with \verb|\hypersetup{...}|.
A clickable TOC comes free once \texttt{hyperref} is loaded.

% =====================================================================
\section{Code Listings}
\begin{lstlisting}[style=code, language=Python, caption={Demo}, label={lst:demo}]
def greet(name):
    """Say hello."""
    return f"Hello, {name}"      # f-string
\end{lstlisting}
Inline code: \verb|\lstinline|print(x)||. From a file:
\verb|\lstinputlisting[style=code]{script.py}|. For prettier output use
\texttt{minted} (needs \verb|-shell-escape| + Pygments).

% =====================================================================
\section{Theorems, Proofs \& Custom Boxes}
\begin{definition}
A statement being defined.
\end{definition}
\begin{theorem}[Named]\label{thm:x}
The claim.
\end{theorem}
\begin{proof}
Argument. \qedhere
\end{proof}
\begin{proposition}
Shares the theorem counter (see the preamble \verb|[theorem]| option).
\end{proposition}
A coloured ``custom theorem box'' with \texttt{tcolorbox}:
\begin{keybox}[Key Result]
$E = mc^2$ inside a coloured, titled box.
\end{keybox}

% =====================================================================
\section{Algorithms / Pseudocode}
\begin{algorithm}[H]
\SetAlgoLined
\KwIn{array $A$ of length $n$}
\KwOut{maximum element}
$m \gets A[1]$\;
\For{$i \gets 2$ \KwTo $n$}{
  \If{$A[i] > m$}{ $m \gets A[i]$\; }
}
\Return $m$\;
\caption{Find maximum (algorithm2e)}
\end{algorithm}
Alternative package: \texttt{algorithm} + \texttt{algpseudocode}
(\verb|\State \For \While \If \EndIf|).

% =====================================================================
\section{Footnotes}
Text with a footnote.\footnote{This is the footnote body.}
Multiple footnotes auto-number. Reuse one with
\verb|\footnotemark| / \verb|\footnotetext|.

% =====================================================================
\section{References / Bibliography}
Manual (no external tool, fastest for exams):
\begin{verbatim}
\begin{thebibliography}{9}
  \bibitem{key1} A. Author, \textit{Title}, Publisher, Year.
  \bibitem{key2} B. Writer, ``Article,'' Journal, 2024.
\end{thebibliography}
\end{verbatim}
Cite in text with \verb|\cite{key1}| $\to$ [1]. The \verb|{9}| argument
just sizes the label box (one digit here).

% =====================================================================
\section{Predicted Tricky Variations}
\begin{enumerate}[label=\arabic*.]
  \item \textbf{Aligned multi-column equations}: use \texttt{alignat} to
        control inter-column spacing manually:
        \verb|\begin{alignat}{2} a&=b &\quad c&=d \end{alignat}|.
  \item \textbf{Piecewise + annotations}: \texttt{cases} plus
        \verb|\underbrace{expr}_{\text{note}}| and
        \verb|\overbrace{expr}^{\text{note}}|.
  \item \textbf{Numbered sub-equations}: wrap several equations in
        \verb|\begin{subequations} ... \end{subequations}| to get
        (1a),(1b).
  \item \textbf{Long division-style derivations}: keep \verb|&=| aligned
        and end intermediate lines with \verb|\nonumber| so only the
        final result is numbered.
  \item \textbf{Units}: load \texttt{siunitx} and use \verb|\SI{9.8}{m/s^2}|
        and \verb|\num{1.496e11}| for clean scientific notation.
  \item \textbf{Coloured full-row headers}: \verb|\rowcolor| must be the
        FIRST token of the row, before any \verb|\multicolumn|.
  \item \textbf{Landscape table/figure}: \texttt{\{sidewaystable\}} or the
        \texttt{lscape} \verb|{landscape}| environment.
  \item \textbf{Custom theorem box with counter}: \texttt{tcolorbox} with
        \verb|\newtcbtheorem| gives numbered, coloured theorems.
  \item \textbf{Chemical/tensor indices}: pre-subscripts via
        \verb|{}_a^b X| and aligned indices via phantom \verb|\ | or
        \verb|\vphantom|.
  \item \textbf{Boxed final answer}: \verb|\boxed{x = \frac{-b\pm\sqrt{\Delta}}{2a}}|.
\end{enumerate}

Referenced items for completeness: Equation~\eqref{eq:cs1},
Theorem~\ref{thm:x}, Figure~\ref{fig:subs}, Listing~\ref{lst:demo}.

\end{document}
